% =============================================================================
%  CHAPITRE 3 — PRÉSENTATION DU SITE DE L'ÉTUDE, DES DONNÉES ET DES RÉSULTATS
% =============================================================================
\entete{Chapitre 3 — Site de l'étude, données et résultats}
\chapter{Présentation du site de l'étude, des données et des résultats}
\label{chap:terrain}

\introchap

Ce chapitre confronte la démarche exposée précédemment au terrain. Il présente
d'abord la Fondation Jean Félicien Gacha et le \emph{Think Tank} Foudjem, cadre
concret de l'investigation (section~\ref{sec:site}), puis le déroulement du
stage (section~\ref{sec:stage}). Il expose ensuite les données recueillies et
l'étude de l'existant (section~\ref{sec:donnees}), dont découle le cahier des
charges et l'évaluation du coût du projet (section~\ref{sec:cdc}). Vient enfin
la partie proprement technique : la modélisation du système
(section~\ref{sec:modelisation}), son implémentation
(section~\ref{sec:implementation}) et la présentation détaillée des résultats
obtenus, fonctionnalité par fonctionnalité (section~\ref{sec:resultats}).

% =============================================================================
\section{Présentation du site de l'étude}
\label{sec:site}

\subsection{Identification de la structure}

La Fondation Jean Félicien Gacha est une organisation à but non lucratif de
droit camerounais, active dans les domaines de la culture, de l'art, de
l'artisanat, de l'éducation et du développement communautaire. Elle est
implantée à Bangoulap, dans l'arrondissement de Bangangté, département du Ndé,
région de l'Ouest du Cameroun. Le tableau~\ref{tab:fiche} en présente la fiche
d'identification.

\begin{table}[htbp]
\centering
\caption{Fiche d'identification de la Fondation Jean Félicien Gacha}
\label{tab:fiche}
\small
\begin{tabularx}{\textwidth}{|L{5.2cm}|X|}
\hline
\rowcolor{keycetete}
\thead{Désignation} & \thead{Affectation} \\ \hline
Dénomination & Fondation Jean Félicien Gacha \\ \hline
Nature juridique & Fondation — organisation à but non lucratif \\ \hline
Secteur d'activité & Culture, art, artisanat, éducation, développement communautaire \\ \hline
Siège & Bangoulap, arrondissement de Bangangté, département du Ndé, région de l'Ouest \\ \hline
Pays & Cameroun \\ \hline
Entité d'accueil du stage & \emph{Think Tank} Foudjem \\ \hline
Encadreur professionnel & M. Adrien GHONSI \\ \hline
Période du stage & Du 20 juin au 20 août 2026 \\ \hline
Domaines d'intervention & Espaces d'exposition, ateliers de création, résidences d'artistes, espaces verts, actions communautaires \\ \hline
\end{tabularx}
\source{Documents internes de la Fondation et entretiens, juin 2026.}
\end{table}

\encadre{Précision à apporter}{%
Les rubriques laissées volontairement génériques ci-dessus — date de création,
effectif du personnel, coordonnées officielles, numéro d'enregistrement — doivent
être complétées à partir des documents administratifs de la Fondation. Nous
avons préféré ne rien avancer que nous n'ayons pu vérifier.}

\subsection{Historique et évolution}

La Fondation porte le nom de son initiateur, Jean Félicien Gacha, dont l'action
s'est inscrite dans une volonté de faire de Bangoulap un lieu de rencontre entre
la création contemporaine et les savoir-faire hérités des chefferies de l'Ouest.
Née d'une initiative privée, la structure s'est progressivement dotée d'espaces
d'exposition, d'ateliers et de résidences, puis a élargi son action au
développement communautaire local. Son évolution suit ainsi un mouvement
caractéristique : d'une collection privée vers une institution ouverte au
public, puis vers un acteur du développement territorial.

Cette trajectoire explique l'un des traits saillants du terrain : la Fondation
dispose d'un patrimoine matériel et documentaire considérable, mais d'une
organisation encore largement fondée sur la mémoire des personnes et sur des
supports non structurés — cahiers, plaquettes, dossiers photographiques
dispersés. C'est très exactement le point d'appui de notre intervention.

\subsection{Missions et activités principales}

Les missions de la Fondation peuvent être regroupées en quatre ensembles :

\begin{itemize}
  \item \textbf{la conservation et la présentation du patrimoine} : réunir,
        documenter et exposer des objets et des témoignages représentatifs de la
        culture de l'Ouest camerounais ;
  \item \textbf{le soutien à la création} : accueillir des artistes en
        résidence, mettre à leur disposition des ateliers et diffuser leur
        production ;
  \item \textbf{la transmission} : accueillir des scolaires, des chercheurs et
        des visiteurs, et faire connaître les savoir-faire artisanaux ;
  \item \textbf{le développement communautaire} : conduire des actions au
        bénéfice des populations environnantes.
\end{itemize}

\subsection{Organisation et fonctionnement}

La Fondation s'organise autour d'une direction, de pôles opérationnels
— conservation et exposition, ateliers et résidences, accueil du public,
communication, administration — et d'une cellule de réflexion et d'innovation,
le \emph{Think Tank} Foudjem, au sein de laquelle notre stage s'est déroulé.

Le \emph{Think Tank} Foudjem a pour vocation d'identifier les leviers de
modernisation de l'institution et d'expérimenter des réponses. C'est de cette
cellule, et plus précisément de son responsable M. Adrien GHONSI, qu'est venue
la formulation initiale du problème traité dans ce rapport : \emph{comment
permettre à quelqu'un qui ne peut pas venir de visiter tout de même la
Fondation ?}

\begin{figure}[htbp]
\centering
\aplacer{7.2cm}{Représenter l'organigramme de la Fondation Jean Félicien Gacha :
la direction générale au sommet ; en dessous, les pôles Conservation et
exposition, Ateliers et résidences, Accueil du public, Communication,
Administration et finances ; et, rattaché à la direction, le \emph{Think Tank}
Foudjem. Mettre ce dernier en évidence (encadré de couleur), puisque c'est
l'entité d'accueil du stage.}
\caption{Organigramme de la Fondation Jean Félicien Gacha}
\label{fig:organigramme}
\source{Élaboré par nos soins à partir des entretiens conduits en juin 2026.}
\end{figure}

\subsection{Environnement technologique existant}

Le relevé de l'environnement numérique de la Fondation, effectué durant la
première semaine du stage, a mis en évidence la situation suivante :

\begin{itemize}
  \item \textbf{aucun système de gestion des collections} : l'inventaire des
        œuvres était tenu sur des supports bureautiques dispersés et sur des
        cahiers ;
  \item \textbf{une présence en ligne réduite} : quelques pages de présentation
        et des publications sur les réseaux sociaux, sans base documentaire ;
  \item \textbf{aucune billetterie ni boutique en ligne} ; les entrées et les
        ventes d'artisanat se traitaient exclusivement sur place ;
  \item \textbf{aucun fichier des visiteurs} : une personne ayant visité la
        Fondation ne pouvait pas être recontactée pour être informée d'une
        exposition ;
  \item \textbf{aucune sauvegarde structurée} des photographies et des notices.
\end{itemize}

\subsection{Justification du choix du terrain}

Le choix de ce terrain se justifie par trois caractéristiques qui en font un cas
particulièrement instructif. D'abord, la richesse patrimoniale y est réelle et
documentée, ce qui exclut le risque de construire un dispositif sans matière.
Ensuite, les contraintes de moyens y sont représentatives de celles que
rencontrent la plupart des institutions patrimoniales de la sous-région : ce qui
fonctionne ici a des chances de fonctionner ailleurs. Enfin, la présence du
\emph{Think Tank} Foudjem garantissait un interlocuteur capable de formuler le
besoin, d'arbitrer les choix et d'évaluer les propositions — condition
indispensable à une recherche interactive.

% =============================================================================
\section{Déroulement du stage}
\label{sec:stage}

\subsection{Accueil et intégration}

Le stage a débuté le 20 juin 2026 par un entretien avec M. Adrien GHONSI,
responsable du \emph{Think Tank} Foudjem. Cet échange avait pour objet de
préciser les attentes de la structure, de nous informer du règlement intérieur
et de définir les objectifs de la période. C'est au cours de cette rencontre que
le problème de la visite immersive nous a été posé.

L'intégration s'est faite en quelques jours, facilitée par la disponibilité du
personnel. Une visite complète des espaces, conduite par le pôle Accueil du
public, nous a permis de prendre la mesure du lieu — condition préalable à toute
tentative de le restituer numériquement. Nous avons pu, dès la première semaine,
observer des visites guidées réelles et relever les questions que les visiteurs
posaient spontanément.

\subsection{Chronogramme des tâches effectuées}

Les travaux ont été supervisés par l'encadreur professionnel selon un programme
hebdomadaire présenté au tableau~\ref{tab:chronogramme}.

\begin{table}[htbp]
\centering
\caption{Chronogramme des tâches effectuées durant le stage}
\label{tab:chronogramme}
\footnotesize
\begin{tabularx}{\textwidth}{|L{3.4cm}|X|}
\hline
\rowcolor{keycetete}
\thead{Période} & \thead{Tâches effectuées} \\ \hline
Du 20 au 26 juin 2026 &
Accueil et présentation de la Fondation ; prise de contact avec le personnel du
\emph{Think Tank} Foudjem ; visite complète des espaces ; formulation initiale
du besoin ; observation de visites guidées. \\ \hline
Du 27 juin au 3 juillet 2026 &
Étude de l'existant ; entretiens semi-directifs avec six membres du personnel ;
relevé de l'environnement numérique ; dépouillement des inventaires et des
notices ; administration du questionnaire aux visiteurs. \\ \hline
Du 4 au 10 juillet 2026 &
Rédaction du cahier des charges ; arbitrage des besoins avec l'encadreur
professionnel ; modélisation UML (cas d'utilisation, classes, séquences) ;
choix de l'architecture et des technologies. \\ \hline
Du 11 au 17 juillet 2026 &
Conception et création de la base de données ; mise en place du cloisonnement
par organisation ; développement du back-office : musées, salles, œuvres,
tarifs, médias. \\ \hline
Du 18 au 24 juillet 2026 &
Développement du site public ; catalogue, fiches d'œuvres, boutique,
billetterie avec billets à code QR ; gestion des événements et du livre d'or. \\ \hline
Du 25 au 31 juillet 2026 &
Développement du moteur de visite immersive à 360~degrés ; éditeur de parcours
dans le back-office ; pose des points d'intérêt ; enchaînement des salles. \\ \hline
Du 1\textsuperscript{er} au 7 août 2026 &
Numérisation en trois dimensions ; réalité augmentée ; génération des codes QR ;
essais de photogrammétrie et d'acquisition par capteur LiDAR ; couverture des
appareils Android et iOS. \\ \hline
Du 8 au 14 août 2026 &
Développement de l'assistant conversationnel à ancrage documentaire ;
rapprochement documentaire avec les collections ouvertes ; module de généalogie ;
campagnes d'information aux adhérents. \\ \hline
Du 15 au 20 août 2026 &
Conteneurisation ; chaînes d'intégration et de déploiement continus ;
description de l'infrastructure ; mise en ligne ; essais de bout en bout ;
mesures de performance et de coût ; restitution à l'encadreur professionnel. \\ \hline
\end{tabularx}
\source{Journal de bord du stage, juin--août 2026.}
\end{table}

\subsection{Acquis du stage}

Cette période a permis de consolider des connaissances théoriques et d'en
acquérir de nouvelles. Nous avons ainsi eu l'occasion :

\begin{itemize}
  \item de conduire un projet complet, de la formulation du besoin à la mise en
        production, dans un cadre professionnel réel ;
  \item de mettre en pratique la modélisation UML et le processus 2TUP sur un
        système de taille significative ;
  \item de concevoir et de mettre en œuvre un cloisonnement des données au
        niveau de la base de données, et d'en éprouver la robustesse ;
  \item de nous approprier les architectures de génération augmentée par la
        recherche et les mécanismes de plongement lexical ;
  \item d'écrire un moteur de rendu tridimensionnel sur mesure, faute de pouvoir
        recourir à une bibliothèque existante ;
  \item de pratiquer l'intégration et le déploiement continus, la
        conteneurisation et la description de l'infrastructure par le code ;
  \item de mesurer et d'interpréter le coût d'une infrastructure en nuage selon
        les principes du FinOps ;
  \item de nous familiariser avec la conduite d'entretiens et l'exploitation
        d'une enquête par questionnaire.
\end{itemize}

\subsection{Difficultés rencontrées}

Trois difficultés méritent d'être signalées, car elles ont infléchi le travail.

\begin{enumerate}
  \item \textbf{L'impossibilité d'installer de nouvelles bibliothèques
        logicielles.} Cette contrainte a imposé d'écrire sur mesure le moteur de
        panorama, le moteur de relief, l'encodeur de codes QR et le générateur
        de fichiers tridimensionnels. Le coût en temps a été élevé ; le bénéfice
        — un système sans dépendance pour ses fonctions les plus visibles — n'a
        été perçu qu'à la fin.

  \item \textbf{L'absence de matière numérique préexistante.} Il n'existait ni
        photographies à 360~degrés, ni modèles tridimensionnels, ni notices
        structurées. Nous avons donc conçu un mode de démonstration capable de
        fonctionner sur une base vierge : une salle et deux objets sont générés
        par le programme lui-même, ce qui permet de présenter le dispositif
        avant même les prises de vue définitives.

  \item \textbf{L'indisponibilité de certaines interfaces d'accès.} Quatre
        grandes collections exigent une clé que nous n'avons pu obtenir, et
        certaines fonctions du navigateur nécessaires aux essais de réalité
        augmentée ne sont pas disponibles sur poste Windows. Ces essais ont dû
        être conduits sur téléphone, sur un nombre restreint d'appareils.
\end{enumerate}

% =============================================================================
\section{Présentation des données collectées}
\label{sec:donnees}

\subsection{Étude de l'existant}

L'étude de l'existant consiste à analyser les outils, les processus et les
systèmes en place afin d'en identifier les forces, les faiblesses et les besoins
non satisfaits. Appliquée à la Fondation, elle révèle un fonctionnement
entièrement manuel pour ce qui touche à la valorisation.

La découverte des espaces suppose une visite sur place. Les collections ne sont
inventoriées que partiellement, sur des supports non reliés entre eux. Les
tarifs et les entrées sont gérés sur registre papier. La communication vers le
public repose sur des publications ponctuelles sur les réseaux sociaux, sans
possibilité de recontacter les personnes intéressées. Enfin, aucune information
n'est disponible pour le visiteur en dehors de la présence d'un guide.

Le tableau~\ref{tab:existant} récapitule les insuffisances relevées et les
besoins qui en découlent.

\begin{table}[htbp]
\centering
\caption{Insuffisances du dispositif existant et besoins induits}
\label{tab:existant}
\small
\begin{tabularx}{\textwidth}{|L{4.4cm}|X|L{3.4cm}|}
\hline
\rowcolor{keycetete}
\thead{Domaine} & \thead{Insuffisance constatée} & \thead{Besoin induit} \\ \hline
Accès au lieu & La visite exige un déplacement physique & Parcours immersif à distance \\ \hline
Perception des œuvres & Photographies plates, aucune notion de volume ni de taille & Présentation en trois dimensions et en réalité augmentée \\ \hline
Documentation & Notices dispersées, objets traités isolément & Inventaire structuré et mise en relation \\ \hline
Information du visiteur & Aucune réponse possible hors présence d'un guide & Assistant conversationnel ancré sur les documents \\ \hline
Gestion administrative & Registres papier, données non sauvegardées & Progiciel de gestion intégré \\ \hline
Relation au public & Aucun fichier de visiteurs & Adhésion en ligne et campagnes d'information \\ \hline
Rayonnement & Dispositif non transposable à d'autres institutions & Architecture ouverte à plusieurs organisations \\ \hline
\end{tabularx}
\source{Étude de l'existant, juin--juillet 2026.}
\end{table}

\subsection{Résultats des entretiens}

L'analyse thématique des six entretiens conduits auprès du personnel a fait
émerger cinq thèmes, présentés ici par ordre de fréquence d'apparition.

\begin{enumerate}
  \item \textbf{L'accessibilité} (mentionnée dans les six entretiens). Le
        déplacement est unanimement désigné comme le premier obstacle. La
        diaspora, les chercheurs étrangers et les établissements scolaires
        éloignés sont explicitement cités.

  \item \textbf{La perception de l'espace} (cinq entretiens). Les personnes
        interrogées insistent sur le fait qu'une photographie ne rend compte ni
        des volumes ni de l'atmosphère des salles. L'expression « il faut être
        là pour comprendre » revient à plusieurs reprises.

  \item \textbf{La mise en relation des objets} (quatre entretiens). Ce thème
        n'était pas anticipé. Il est apparu lorsqu'un interlocuteur a évoqué des
        pièces comparables conservées à l'étranger, dont la Fondation n'a
        connaissance que par ouï-dire. Cette remarque a directement engendré la
        deuxième question spécifique de l'étude.

  \item \textbf{La fiabilité de l'information} (quatre entretiens). La crainte
        exprimée est nette : un dispositif automatique qui énoncerait des
        informations fausses sur le patrimoine porterait atteinte à la
        crédibilité de l'institution. Cette exigence a orienté toute la
        conception de l'assistant.

  \item \textbf{La pérennité économique} (trois entretiens). La question posée
        est celle du coût de fonctionnement une fois le stagiaire parti. Elle
        justifie à elle seule la place accordée au FinOps dans ce travail.
\end{enumerate}

\subsection{Résultats du questionnaire}

Le questionnaire a été administré à 42 visiteurs. Le
tableau~\ref{tab:profil} présente le profil des répondants et le
tableau~\ref{tab:attentes} leurs attentes.

\begin{table}[htbp]
\centering
\caption{Profil des répondants au questionnaire}
\label{tab:profil}
\small
\begin{tabularx}{\textwidth}{|L{4.0cm}|X|C{2.2cm}|C{2.2cm}|}
\hline
\rowcolor{keycetete}
\thead{Caractéristique} & \thead{Modalité} & \thead{Effectif} & \thead{Fréquence} \\ \hline
\multirow{3}{*}{Tranche d'âge}
 & 18 à 30 ans & 19 & 45,2 \% \\ \cline{2-4}
& 31 à 45 ans & 15 & 35,7 \% \\ \cline{2-4}
& 46 ans et plus & 08 & 19,1 \% \\ \hline
\multirow{3}{*}{Provenance}
 & Région de l'Ouest & 21 & 50,0 \% \\ \cline{2-4}
& Autre région du Cameroun & 14 & 33,3 \% \\ \cline{2-4}
& Étranger & 07 & 16,7 \% \\ \hline
\multirow{2}{*}{Équipement principal}
 & Téléphone mobile & 36 & 85,7 \% \\ \cline{2-4}
& Ordinateur & 06 & 14,3 \% \\ \hline
\rowcolor{keycetotal}
\textbf{Total} & & \textbf{42} & \textbf{100 \%} \\ \hline
\end{tabularx}
\source{Enquête par questionnaire, juillet 2026.}
\end{table}

Deux enseignements ressortent de ce profil. D'une part, la moitié des visiteurs
vient déjà de l'extérieur de la région, ce qui atteste l'existence d'un public
prêt à s'intéresser à la Fondation à distance. D'autre part, plus de huit
visiteurs sur dix accèdent au numérique par un téléphone mobile : le dispositif
devait donc être conçu pour le téléphone en priorité, et non pour l'ordinateur.

\begin{table}[htbp]
\centering
\caption{Attentes exprimées à l'égard d'un dispositif de visite à distance}
\label{tab:attentes}
\small
\begin{tabularx}{\textwidth}{|X|C{2.4cm}|C{2.4cm}|}
\hline
\rowcolor{keycetete}
\thead{Attente exprimée} & \thead{Effectif} & \thead{Fréquence} \\ \hline
Se déplacer dans les salles comme si l'on y était & 38 & 90,5 \% \\ \hline
Voir les objets sous tous leurs angles et connaître leur taille réelle & 35 & 83,3 \% \\ \hline
Poser des questions et obtenir une réponse immédiate & 33 & 78,6 \% \\ \hline
Savoir à qui l'objet a appartenu et d'où il vient & 29 & 69,0 \% \\ \hline
Savoir si des objets semblables existent ailleurs dans le monde & 24 & 57,1 \% \\ \hline
Acheter un billet ou un objet d'artisanat en ligne & 18 & 42,9 \% \\ \hline
\end{tabularx}
\source{Enquête par questionnaire, juillet 2026 — questions à réponses multiples.}
\end{table}

\begin{figure}[htbp]
\centering
\aplacer{6.5cm}{Insérer un diagramme en barres horizontales reprenant les six
attentes du tableau~\ref{tab:attentes}, classées par fréquence décroissante, avec
les pourcentages en bout de barre. Employer la couleur bleue du rapport pour les
barres et l'or pour la barre la plus élevée.}
\caption{Attentes des visiteurs à l'égard d'un dispositif de visite à distance}
\label{fig:attentes}
\source{Enquête par questionnaire, juillet 2026 ($n = 42$).}
\end{figure}

Ces résultats confirment la hiérarchie retenue dans les objectifs : la
restitution du déplacement dans l'espace arrive en tête, suivie de la perception
des objets, puis de la possibilité d'interroger. Ils valident également la
pertinence de la mise en relation entre institutions, souhaitée par plus d'un
visiteur sur deux alors même qu'elle ne figurait pas dans la demande initiale.

% =============================================================================
\section{Cahier des charges du système}
\label{sec:cdc}

Le cahier des charges consigne les spécifications fonctionnelles et techniques
de l'application. Il sert de référence tout au long du développement et
garantit que le produit livré correspond aux besoins exprimés. La plateforme
construite a été nommée \textbf{MUSÉA}.

\subsection{Besoins fonctionnels}

Les besoins fonctionnels sont regroupés en six familles.

\begin{enumerate}
  \item \textbf{Gestion de l'institution.} Enregistrer et publier les musées,
        les salles et les œuvres ; gérer les médias associés ; définir une
        grille tarifaire ; administrer les événements.

  \item \textbf{Visite immersive.} Composer des parcours à 360~degrés reliant
        les salles ; poser des points d'intérêt renvoyant à des œuvres ;
        permettre au visiteur de se déplacer, de pivoter et de zoomer.

  \item \textbf{Présentation des œuvres.} Afficher chaque œuvre en trois
        dimensions ; permettre sa projection à taille réelle dans
        l'environnement du visiteur ; restituer sa notice, son récit et sa voix.

  \item \textbf{Assistance au visiteur.} Répondre aux questions à partir des
        seuls documents validés ; refuser les questions hors périmètre ;
        consigner les questions restées sans réponse.

  \item \textbf{Mise en relation documentaire.} Rechercher les œuvres
        apparentées dans les collections ouvertes ; les présenter avec leur
        justification ; les soumettre à validation.

  \item \textbf{Relation au public et commerce.} Permettre l'adhésion en ligne ;
        vendre billets et produits ; envoyer des campagnes d'information ;
        recueillir les avis.
\end{enumerate}

\subsection{Besoins non fonctionnels}

\begin{itemize}
  \item \textbf{Performance.} Affichage d'une salle en moins de trois secondes
        sur une connexion mobile ordinaire ; réponse de l'assistant en moins de
        cinq secondes.
  \item \textbf{Sécurité.} Authentification obligatoire pour toute écriture ;
        cloisonnement strict des données entre organisations, appliqué par la
        base de données elle-même ; aucune clé d'accès à un service tiers
        présente dans le navigateur.
  \item \textbf{Robustesse.} Une fonctionnalité dépendant d'un service externe
        indisponible doit se dégrader proprement, sans interrompre le parcours
        du visiteur.
  \item \textbf{Convivialité.} Utilisation sans formation préalable ; interface
        conçue pour le téléphone en priorité ; bilinguisme français-anglais.
  \item \textbf{Accessibilité.} Navigation au clavier ; éléments interactifs
        restituables par les lecteurs d'écran.
  \item \textbf{Portabilité.} Fonctionnement dans un navigateur, sans
        installation ; application installable sur l'écran d'accueil.
\end{itemize}

\subsection{Contraintes technologiques}

Les principales technologies imposées ou retenues ont été présentées au
tableau~\ref{tab:outils}. Trois contraintes structurantes méritent d'être
rappelées ici : l'impossibilité d'ajouter des bibliothèques, l'exigence d'un
fonctionnement sur téléphone d'entrée de gamme, et l'obligation d'un
cloisonnement garanti par la base de données et non par le code applicatif.

\subsection{Contraintes financières : ressources et coût du projet}

\subsubsection{Ressources matérielles}

Le tableau~\ref{tab:materiel} recense le matériel physique mobilisé pour la
réalisation du projet.

\begin{table}[htbp]
\centering
\caption{Ressources matérielles mobilisées}
\label{tab:materiel}
\small
\begin{tabularx}{\textwidth}{|L{3.2cm}|X|C{1.3cm}|R{2.0cm}|R{2.0cm}|}
\hline
\rowcolor{keycetete}
\thead{Matériel} & \thead{Caractéristiques} & \thead{Qté} & \thead{P.U. (FCFA)} & \thead{Total (FCFA)} \\ \hline
Ordinateur portable & Processeur 8 cœurs, 32 Go de mémoire vive, 1 To de disque à semi-conducteurs & 1 & 450 000 & 450 000 \\ \hline
Écran secondaire & 24 pouces, 1920 $\times$ 1080 & 1 & 100 000 & 100 000 \\ \hline
Téléphone à capteur LiDAR & iPhone 15 Pro — acquisition tridimensionnelle et essais de réalité augmentée & 1 & 400 000 & 400 000 \\ \hline
Disque externe & 1 To, à semi-conducteurs — sauvegarde des prises de vue et des maillages & 1 & 55 000 & 55 000 \\ \hline
Onduleur & 650 VA — protection contre les coupures de courant & 1 & 48 000 & 48 000 \\ \hline
Éclairage d'appoint & 2 panneaux à diodes — prises de vue des œuvres & 1 & 45 000 & 45 000 \\ \hline
Trépied et tête panoramique & Prises de vue à 360 degrés & 1 & 35 000 & 35 000 \\ \hline
\rowcolor{keycetotal}
\multicolumn{4}{|l|}{\textbf{Coût total des ressources matérielles}} & \textbf{1 133 000} \\ \hline
\end{tabularx}
\source{Relevé des prix pratiqués à Bafoussam et à Douala, juin 2026.}
\end{table}

\subsubsection{Ressources logicielles et services}

Le tableau~\ref{tab:logiciel} recense les logiciels et services utilisés. Il
convient de souligner que la quasi-totalité de la chaîne technique repose sur
des logiciels libres ou sur des paliers gratuits, ce qui constitue en soi un
résultat : le coût d'entrée d'un tel dispositif n'est pas logiciel, il est
matériel et humain.

\begin{table}[htbp]
\centering
\caption{Ressources logicielles et services mobilisés}
\label{tab:logiciel}
\small
\begin{tabularx}{\textwidth}{|L{4.4cm}|X|R{2.4cm}|}
\hline
\rowcolor{keycetete}
\thead{Logiciel ou service} & \thead{Rôle} & \thead{Coût (FCFA)} \\ \hline
Windows 11 Professionnel & Système d'exploitation du poste de développement & 90 350 \\ \hline
Connexion Internet & Forfait données, deux mois & 60 000 \\ \hline
Nom de domaine \texttt{.store} & Adresse publique de la plateforme, un an & 9 000 \\ \hline
Hébergement et diffusion AWS & Stockage, réseau de diffusion, noms de domaine et certificats, deux mois & 1 460 \\ \hline
Calcul graphique ponctuel & Instance à processeur graphique pour la photogrammétrie, 6 heures & 1 900 \\ \hline
Supabase & Base de données, authentification et fonctions serveur — palier gratuit & 0 \\ \hline
Groq & Génération de texte — palier gratuit & 0 \\ \hline
ElevenLabs & Synthèse vocale — palier de découverte & 0 \\ \hline
Visual Studio Code, Git, Docker, Terraform & Développement et industrialisation — logiciels libres & 0 \\ \hline
Meshroom, Blender & Photogrammétrie et préparation des maillages — logiciels libres & 0 \\ \hline
\rowcolor{keycetotal}
\multicolumn{2}{|l|}{\textbf{Coût total des ressources logicielles et services}} & \textbf{162 710} \\ \hline
\end{tabularx}
\source{Factures et relevés de consommation, juin--août 2026.}
\end{table}

\subsubsection{Coût de l'infrastructure en nuage : le relevé FinOps}

Le coût de l'infrastructure ne se devine pas, il se relève. Conformément aux
principes exposés au chapitre~1, la dépense a été suivie au moyen de l'outil
d'analyse des coûts du fournisseur. Le tableau~\ref{tab:finops} détaille la
dépense par service.

\begin{table}[htbp]
\centering
\caption{Détail du coût de l'infrastructure en nuage (relevé FinOps)}
\label{tab:finops}
\small
\begin{tabularx}{\textwidth}{|L{3.8cm}|X|R{2.2cm}|R{2.2cm}|}
\hline
\rowcolor{keycetete}
\thead{Service} & \thead{Usage} & \thead{Coût mensuel (FCFA)} & \thead{Sur 2 mois (FCFA)} \\ \hline
Route 53 & Zone hébergée du domaine et enregistrements joker & 310 & 620 \\ \hline
Amazon S3 & Stockage du site, des modèles 3D et des médias & 90 & 180 \\ \hline
CloudFront & Réseau de diffusion et certificat de sécurité & 330 & 660 \\ \hline
Certificate Manager & Certificats de sécurité — service gratuit & 0 & 0 \\ \hline
\rowcolor{keycetotal}
\multicolumn{2}{|l|}{\textbf{Total hébergement récurrent}} & \textbf{730} & \textbf{1 460} \\ \hline
EC2 (processeur graphique) & Reconstruction photogrammétrique, 6 heures, instance arrêtée hors usage & — & 1 900 \\ \hline
\rowcolor{keycetotal}
\multicolumn{3}{|l|}{\textbf{Total infrastructure sur la période du stage}} & \textbf{3 360} \\ \hline
\end{tabularx}
\source{AWS Cost Explorer, période du 20 juin au 20 août 2026 (voir
figure~\ref{fig:finops}).}
\end{table}

\begin{figure}[htbp]
\centering
\acapturer{7cm}{Insérer la capture d'écran d'AWS Cost Explorer montrant la
dépense mensuelle ventilée par service sur la période du stage. Veiller à ce que
les montants et les noms de services soient lisibles, et à masquer tout
identifiant de compte.}
\caption{Relevé de la dépense mensuelle d'infrastructure par service}
\label{fig:finops}
\source{Console AWS Cost Explorer, août 2026.}
\end{figure}

Ce relevé appelle un commentaire qui dépasse la simple comptabilité. Le coût
récurrent de l'hébergement s'établit à environ \textbf{730 FCFA par mois}, soit
moins de 9 000 FCFA par an. Ce montant, dérisoire au regard du budget d'une
institution, tient à un choix d'architecture : l'application étant un site
statique servi par un réseau de diffusion, elle ne mobilise aucun serveur en
fonctionnement permanent. Le seul poste réellement coûteux — le calcul graphique
nécessaire à la photogrammétrie — est ponctuel et l'instance correspondante
s'arrête automatiquement lorsqu'elle n'est pas sollicitée. C'est là l'apport
concret de la démarche FinOps : non pas économiser après coup, mais concevoir de
telle sorte que la dépense reste proportionnée à l'usage réel.

\subsubsection{Coût total du projet}

\begin{table}[htbp]
\centering
\caption{Coût total du projet}
\label{tab:couttotal}
\small
\begin{tabularx}{\textwidth}{|X|R{3.4cm}|}
\hline
\rowcolor{keycetete}
\thead{Libellé} & \thead{Montant (FCFA)} \\ \hline
Ressources matérielles & 1 133 000 \\ \hline
Ressources logicielles et services & 162 710 \\ \hline
Sous-total & 1 295 710 \\ \hline
Marge d'erreur (3 \%) & 38 871 \\ \hline
\rowcolor{keycetotal}
\textbf{Coût total du projet} & \textbf{1 334 581} \\ \hline
\end{tabularx}
\source{Élaboré par nos soins à partir des tableaux~\ref{tab:materiel}
et~\ref{tab:logiciel}.}
\end{table}

\subsection{Documents livrables}

Les documents remis à la Fondation à l'issue du projet sont les suivants :

\begin{itemize}
  \item le cahier des charges du projet ;
  \item le dossier de conception (diagrammes UML et modèle de données) ;
  \item le manuel d'utilisation du back-office ;
  \item la procédure de mise en production et de reprise ;
  \item le code source et la description de l'infrastructure, versionnés.
\end{itemize}

% =============================================================================
\section{Modélisation du système}
\label{sec:modelisation}

\subsection{Diagramme de cas d'utilisation}

Le diagramme de cas d'utilisation décrit les interactions entre les acteurs et
le système. Quatre acteurs ont été identifiés.

\begin{itemize}
  \item \textbf{Le visiteur} — public, éventuellement non authentifié. Il
        consulte le catalogue, effectue la visite immersive, examine les œuvres
        en trois dimensions et en réalité augmentée, interroge l'assistant,
        consulte la généalogie, adhère, achète un billet ou un produit.

  \item \textbf{Le personnel de l'institution} (administrateur d'organisation).
        Il gère les musées, les salles, les œuvres, les médias, les tarifs, les
        événements, les parcours immersifs, la boutique, les commandes, les
        adhérents et les campagnes ; il valide les rapprochements documentaires
        et contrôle les billets.

  \item \textbf{Le super-administrateur de la plateforme.} Il approuve les
        demandes d'inscription des nouvelles organisations et supervise
        l'ensemble.

  \item \textbf{Les services externes} (acteurs secondaires) : les collections
        ouvertes interrogées, le service de génération de texte, le service de
        synthèse vocale et le service de paiement.
\end{itemize}

\begin{figure}[htbp]
\centering
\aplacer{8.5cm}{Représenter le diagramme de cas d'utilisation général de la
plateforme MUSÉA. Placer à gauche le Visiteur, à droite le Personnel de
l'institution et le Super-administrateur, et en bas les Services externes.
Faire figurer les paquets de cas d'utilisation : \emph{Consulter le catalogue},
\emph{Effectuer la visite immersive}, \emph{Examiner une œuvre en 3D et en RA},
\emph{Interroger l'assistant}, \emph{Explorer la généalogie}, \emph{Adhérer},
\emph{Acheter un billet ou un produit}, \emph{Gérer les collections},
\emph{Composer un parcours}, \emph{Valider les rapprochements},
\emph{Contrôler un billet}, \emph{Envoyer une campagne}, \emph{Approuver une
organisation}. Marquer les relations \guillemotleft~include~\guillemotright{}
(par exemple : Acheter $\rightarrow$ S'authentifier) et
\guillemotleft~extend~\guillemotright{} (Examiner en 3D $\rightarrow$ Afficher
en réalité augmentée).}
\caption{Diagramme de cas d'utilisation général de la plateforme}
\label{fig:usecase}
\source{Élaboré par nos soins.}
\end{figure}

\subsection{Diagramme de classes}

Le diagramme de classes présente la structure statique du système : les entités,
leurs attributs et leurs relations. Les classes principales sont les suivantes.

\begin{itemize}
  \item \textbf{Organisation} : l'institution abonnée à la plateforme. Elle
        porte son identité, son adresse, son plan et ses quotas. Toute donnée
        métier lui est rattachée.
  \item \textbf{Musée}, \textbf{Salle}, \textbf{Œuvre} : la hiérarchie de
        présentation. Une œuvre appartient à une salle, qui appartient à un
        musée, qui appartient à une organisation.
  \item \textbf{Parcours}, \textbf{Scène}, \textbf{PointInteret} : la visite
        immersive. Un parcours enchaîne des scènes ; une scène porte des points
        d'intérêt renvoyant à une œuvre ou à une autre scène.
  \item \textbf{Personnage}, \textbf{LienGenealogique} : la généalogie. Un
        personnage peut avoir possédé des œuvres et être relié à d'autres
        personnages.
  \item \textbf{ObjetExterne}, \textbf{Rapprochement} : la mise en relation
        documentaire. Un objet externe provient d'une collection ouverte ; un
        rapprochement relie une œuvre locale à un objet externe et porte un
        score, une justification et un statut de validation.
  \item \textbf{Profil}, \textbf{Adhesion} : les personnes. Un profil est un
        compte ; une adhésion rattache un visiteur à une organisation.
  \item \textbf{Produit}, \textbf{Commande}, \textbf{LigneCommande},
        \textbf{Acces} : le commerce et la billetterie.
\end{itemize}

\begin{figure}[htbp]
\centering
\aplacer{9cm}{Représenter le diagramme de classes du système. Faire apparaître
la classe \textbf{Organisation} comme racine du cloisonnement, reliée par une
composition à \textbf{Musée} (1..n), lui-même à \textbf{Salle} (1..n), lui-même à
\textbf{Œuvre} (1..n). Rattacher \textbf{Parcours} à \textbf{Musée},
\textbf{Scène} à \textbf{Parcours} et \textbf{PointInteret} à \textbf{Scène},
avec une association de \textbf{PointInteret} vers \textbf{Œuvre}. Faire figurer
\textbf{Personnage} avec l'auto-association \emph{LienGenealogique} et
l'association n..n vers \textbf{Œuvre}. Ajouter \textbf{ObjetExterne} relié à
\textbf{Œuvre} par la classe-association \textbf{Rapprochement} portant les
attributs \emph{score}, \emph{justification} et \emph{statut}. Enfin,
\textbf{Profil}, \textbf{Adhesion}, \textbf{Produit}, \textbf{Commande},
\textbf{LigneCommande} et \textbf{Acces}. Indiquer les attributs principaux et
les cardinalités.}
\caption{Diagramme de classes du système}
\label{fig:classes}
\source{Élaboré par nos soins.}
\end{figure}

\subsection{Diagrammes de séquence}

Deux scénarios ont été retenus pour illustrer la dynamique du système : la
consultation d'une œuvre depuis la visite immersive, et l'interrogation de
l'assistant conversationnel.

\begin{figure}[htbp]
\centering
\aplacer{8cm}{Représenter le diagramme de séquence du scénario
\emph{Consulter une œuvre depuis la visite immersive}. Acteurs et participants :
Visiteur, Interface de visite, Moteur de panorama, Service de données, Base de
données, Visionneuse 3D. Enchaînement : le Visiteur ouvre le parcours ;
l'Interface demande la scène au Service de données, qui l'obtient de la Base ;
le Moteur affiche le panorama et positionne les points d'intérêt ; le Visiteur
touche un point ; l'Interface demande la fiche de l'œuvre ; la fiche s'affiche ;
le Visiteur demande la vue en trois dimensions ; la Visionneuse charge le modèle
et, si l'appareil le permet, propose la réalité augmentée.}
\caption{Diagramme de séquence : consultation d'une œuvre depuis la visite immersive}
\label{fig:seq1}
\source{Élaboré par nos soins.}
\end{figure}

\begin{figure}[htbp]
\centering
\aplacer{8cm}{Représenter le diagramme de séquence du scénario
\emph{Interroger l'assistant conversationnel}. Participants : Visiteur,
Interface de discussion, Fonction serveur \emph{guide-agent}, Base documentaire
vectorielle, Collections ouvertes, Modèle de langue. Enchaînement : le Visiteur
pose une question ; la Fonction serveur applique les garde-fous (question hors
périmètre ? salutation ?) ; elle convertit la question en vecteur ; elle
interroge la Base documentaire et obtient les fragments pertinents ; le cas
échéant elle interroge les Collections ouvertes ; elle compose l'instruction et
appelle le Modèle de langue ; celui-ci rédige la réponse ancrée ; la réponse et
ses sources sont renvoyées au Visiteur ; la question est journalisée de manière
anonyme.}
\caption{Diagramme de séquence : interrogation de l'assistant conversationnel}
\label{fig:seq2}
\source{Élaboré par nos soins.}
\end{figure}

\subsection{Modèle physique des données}

La base de données compte 38 tables, réparties en huit domaines fonctionnels
présentés au tableau~\ref{tab:tables}. Toutes les tables métier portent une
référence à l'organisation propriétaire et sont protégées par une règle de
sécurité au niveau de la ligne.

\begin{table}[htbp]
\centering
\caption{Répartition des tables de la base de données par domaine fonctionnel}
\label{tab:tables}
\small
\begin{tabularx}{\textwidth}{|L{4.2cm}|C{1.6cm}|X|}
\hline
\rowcolor{keycetete}
\thead{Domaine} & \thead{Tables} & \thead{Principales tables} \\ \hline
Organisations et comptes & 6 & \texttt{tenants}, \texttt{tenant\_members}, \texttt{tenant\_partnerships}, \texttt{slugs\_reserves}, \texttt{profiles}, \texttt{site\_settings} \\ \hline
Collections & 6 & \texttt{museums}, \texttt{sectors}, \texttt{objects}, \texttt{object\_tariffs}, \texttt{object\_rarity}, \texttt{object\_views} \\ \hline
Mise en relation & 2 & \texttt{object\_siblings}, \texttt{objets\_externes} \\ \hline
Généalogie & 4 & \texttt{personnages}, \texttt{genealogy\_links}, \texttt{migrations\_historiques}, \texttt{object\_personnage} \\ \hline
Assistance et voix & 5 & \texttt{faq}, \texttt{guide\_questions}, \texttt{voice\_assistants}, \texttt{audio\_tracks}, \texttt{ai\_agent\_config} \\ \hline
Visite immersive et jeux & 6 & \texttt{tours}, \texttt{tour\_scenes}, \texttt{scene\_hotspots}, \texttt{quests}, \texttt{quest\_steps}, \texttt{quest\_progress} \\ \hline
Commerce et accès & 7 & \texttt{products}, \texttt{orders}, \texttt{order\_items}, \texttt{user\_access}, \texttt{subscription\_plans}, \texttt{donation\_tiers}, \texttt{events} \\ \hline
Relation au public & 3 & \texttt{reviews}, \texttt{campaigns}, \texttt{email\_log} \\ \hline
\rowcolor{keycetotal}
\textbf{Total} & \textbf{38} & \\ \hline
\end{tabularx}
\source{Relevé du schéma de la base de données, août 2026.}
\end{table}

% =============================================================================
\section{Implémentation}
\label{sec:implementation}

\subsection{Architecture générale de la plateforme}

L'architecture retenue comporte trois couches. Une \textbf{couche de
présentation} s'exécute dans le navigateur du visiteur : elle regroupe le site
public, le back-office et les moteurs de rendu écrits sur mesure. Une
\textbf{couche de services} regroupe la base de données, l'authentification et
douze fonctions serveur qui exécutent tout ce qui exige une clé d'accès ou un
privilège. Une \textbf{couche de diffusion} sert les fichiers statiques depuis
un réseau de distribution mondial.

Le principe directeur est simple et n'a jamais été enfreint : \emph{aucune clé
d'accès à un service tiers ne se trouve dans le navigateur}. Toute demande
adressée à un modèle de langue, à un service de synthèse vocale ou à un service
de paiement transite par une fonction serveur qui détient seule le secret.

\begin{figure}[htbp]
\centering
\aplacer{9cm}{Représenter l'architecture générale en trois couches.
\textbf{Couche cliente} (navigateur) : Site public, Back-office, Moteur de
panorama WebGL, Visionneuse 3D/RA, Moteur de relief, Encodeur de codes QR.
\textbf{Couche de services} : Base de données PostgreSQL avec sécurité au niveau
de la ligne, Service d'authentification, Stockage de fichiers, et les fonctions
serveur (\emph{guide-agent}, \emph{memory-search}, \emph{freres},
\emph{object-ai}, \emph{campaign-ai}, \emph{setup-agent}, \emph{tts},
\emph{send-email}, \emph{payment-create}, \emph{payment-webhook},
\emph{verifier-domaine}, \emph{europeana}). \textbf{Couche de diffusion} :
stockage d'objets et réseau de distribution. \textbf{Services externes} :
collections ouvertes, service de génération de texte, service de synthèse
vocale, service de paiement. Faire figurer clairement que les clés d'accès ne
franchissent jamais la frontière entre la couche de services et la couche
cliente.}
\caption{Architecture générale de la plateforme MUSÉA}
\label{fig:archi}
\source{Élaboré par nos soins.}
\end{figure}

\subsection{Architecture de l'assistant conversationnel à ancrage documentaire}

L'assistant repose sur une architecture de génération augmentée par la
recherche, dont le principe a été exposé au chapitre~1. Sa mise en œuvre suit
cinq temps.

\begin{enumerate}
  \item \textbf{Ingestion.} Les documents de référence de la Fondation —
        notices d'œuvres, descriptions de salles, histoires des musées,
        biographies des personnages, foire aux questions — sont exportés depuis
        la base de données et déposés dans un espace de stockage d'objets. Ce
        dépôt constitue la source de vérité documentaire : rien de ce qui n'y
        figure pas ne peut fonder une réponse.

  \item \textbf{Découpage et vectorisation.} Chaque document est découpé en
        fragments de taille homogène, puis converti en vecteur numérique par un
        modèle de plongement. Un point de méthode s'est révélé décisif : les
        fragments sont plongés \textbf{en anglais}. Les mesures conduites
        montrent en effet qu'une requête française classait une peinture sans
        rapport devant deux masques réellement apparentés, tandis qu'en anglais
        les correspondances légitimes atteignaient un score d'au moins 0,898
        contre au plus 0,779 pour le bruit — une étendue de 0,20 contre 0,13.

  \item \textbf{Indexation.} Les vecteurs sont enregistrés dans la base de
        données, dans une colonne vectorielle de dimension 384 dotée d'un index
        de type HNSW en distance cosinus. Cet index rend la recherche des plus
        proches voisins praticable sans parcourir l'ensemble des fragments.

  \item \textbf{Recherche et augmentation.} À la réception d'une question, la
        fonction serveur applique d'abord ses garde-fous : elle répond avec
        courtoisie aux salutations, refuse poliment les sujets hors périmètre,
        et restreint le cas échéant la recherche au musée ou à la salle où se
        trouve le visiteur. Elle convertit ensuite la question en vecteur,
        récupère les fragments les plus proches et compose l'instruction
        adressée au modèle de langue.

  \item \textbf{Génération contrôlée.} Le modèle rédige la réponse. La consigne
        lui interdit d'énoncer un fait absent des fragments fournis, lui impose
        de citer l'institution et le numéro d'inventaire lorsqu'ils existent, et
        l'autorise explicitement à déclarer qu'aucune conclusion n'est possible.
        Cette autorisation est essentielle : sans elle, un modèle de langue
        préfère inventer plutôt que de décevoir.
\end{enumerate}

\begin{figure}[htbp]
\centering
\aplacer{9.5cm}{Représenter l'architecture de l'assistant conversationnel à
ancrage documentaire (RAG). \textbf{Chaîne d'ingestion} (partie haute, hors
ligne) : Base de données des notices $\rightarrow$ Export $\rightarrow$
\textbf{Amazon S3} (dépôt documentaire) $\rightarrow$ Découpage en fragments
$\rightarrow$ Traduction et normalisation en anglais $\rightarrow$ \textbf{Modèle
de plongement \texttt{gte-small}} $\rightarrow$ vecteurs de dimension 384
$\rightarrow$ \textbf{Index vectoriel HNSW} (\texttt{pgvector}).
\textbf{Chaîne d'interrogation} (partie basse, en ligne) : Visiteur
$\rightarrow$ Question $\rightarrow$ \textbf{Fonction serveur guide-agent}
$\rightarrow$ garde-fous (hors périmètre, salutation, périmètre musée/salle)
$\rightarrow$ vectorisation de la question $\rightarrow$ recherche des $k$ plus
proches voisins dans l'index $\rightarrow$ fragments retenus $\rightarrow$
composition de l'instruction avec consigne d'ancrage $\rightarrow$
\textbf{Modèle de langue} $\rightarrow$ Réponse sourcée $\rightarrow$ Visiteur.
Ajouter, en dérivation, la branche facultative \emph{Collections ouvertes} et la
branche \emph{Journalisation anonyme des questions sans réponse}. Faire
clairement apparaître que la clé d'accès au modèle de langue reste du côté
serveur.}
\caption{Architecture de l'assistant conversationnel à ancrage documentaire (RAG)}
\label{fig:rag}
\source{Élaboré par nos soins.}
\end{figure}

Une contrainte d'exploitation mérite d'être signalée, car elle a dicté un choix
d'implémentation. La plateforme d'exécution des fonctions serveur impose une
limite de ressources : la vectorisation de dix-huit fragments passe en
1,7~seconde, tandis que vingt fragments déclenchent l'interruption de la
fonction. Les traitements sont donc découpés en lots de douze, valeur retenue
avec une marge de sécurité délibérée.

\subsection{Le moteur de visite immersive}

L'impossibilité d'ajouter des bibliothèques a conduit à écrire le moteur de
panorama de bout en bout, directement au moyen de l'interface graphique du
navigateur. Trois décisions techniques en résument la conception.

\begin{itemize}
  \item \textbf{Un quadrilatère plein écran plutôt qu'une sphère maillée.} Le
        rendu ne projette pas l'image sur une sphère composée de facettes, mais
        calcule pour chaque pixel le rayon correspondant. Cette approche
        supprime les artefacts de facettes visibles sur les maillages grossiers.

  \item \textbf{La mise à la puissance de deux des images.} Les images sont
        redimensionnées afin d'autoriser le mode de répétition de texture, ce
        qui efface la couture verticale visible à la jonction des bords gauche
        et droit du panorama.

  \item \textbf{Une convention d'orientation unique.} L'angle horizontal vaut
        0~degré au centre de l'image et croît vers la droite ; l'angle vertical
        vaut 0~degré à l'horizon et croît vers le haut. Cette convention est
        partagée par la caméra et par les points d'intérêt : sans elle, les
        pastilles se déplacent dans le sens opposé à la caméra.
\end{itemize}

Les points d'intérêt ne sont pas dessinés dans l'image : ce sont de véritables
boutons positionnés en pourcentage au-dessus du rendu. Ce choix, apparemment
mineur, conditionne l'accessibilité : les points restent atteignables au clavier
et restituables par un lecteur d'écran.

\subsection{Chaîne de numérisation tridimensionnelle et réalité augmentée}

Deux chaînes d'acquisition coexistent. La première repose sur la
\textbf{photogrammétrie} : trois orbites de prises de vue guidées autour de
l'objet, puis reconstruction sur une instance à processeur graphique. La
seconde exploite le \textbf{capteur LiDAR} du téléphone, qui fournit directement
une géométrie correctement dimensionnée.

La restitution en réalité augmentée a exigé un traitement particulier de la
diversité des appareils. Trois colonnes ont été ajoutées à la description de
chaque œuvre :

\begin{itemize}
  \item un \textbf{modèle au format USDZ} : les appareils Apple ne déclenchent
        de session de réalité augmentée qu'à partir de ce format, quel que soit
        l'hébergement ;
  \item un \textbf{mode de pose} — au sol ou au mur : sans cette indication, un
        masque destiné à être accroché se retrouve couché sur le plancher ;
  \item un \textbf{facteur d'échelle} : un maillage exporté en centimètres
        arriverait cent fois trop grand.
\end{itemize}

L'objet est présenté à ses \textbf{dimensions réelles}, mesurées sur le maillage
et non saisies à la main. Ce point n'est pas cosmétique : autoriser
l'agrandissement libre ruinerait le propos, qui est de faire éprouver la taille
véritable de la pièce. Le modèle de démonstration produit pour les essais mesure
ainsi 45 $\times$ 52 $\times$ 45 centimètres pour 835 sommets.

Enfin, la réalité augmentée suppose un téléphone, alors qu'une soutenance se
tient devant un ordinateur. Un encodeur de codes QR a donc été écrit afin que le
public puisse scanner l'écran et faire apparaître l'objet dans la salle où il se
trouve. La conformité de l'encodeur a été vérifiée par le décompte exact des
modules de service imposé par la norme et par l'annulation des syndromes du code
correcteur.

\subsection{Infrastructure d'hébergement}

L'application étant un site statique, l'infrastructure retenue est délibérément
minimale : un espace de stockage privé, chiffré et versionné, servi par un
réseau de diffusion mondial au moyen d'un contrôle d'accès à l'origine ; un
certificat de sécurité couvrant le domaine principal et l'ensemble de ses
sous-domaines ; et des enregistrements de noms de domaine dotés d'un caractère
générique, qui font exister l'adresse de chaque organisation sans qu'il soit
nécessaire de créer un enregistrement par institution.

Deux réglages conditionnent le bon fonctionnement du dispositif :

\begin{itemize}
  \item \textbf{le repli applicatif} : un stockage privé répond par un refus
        d'accès, et non par une page absente, lorsqu'une adresse profonde est
        demandée ; les deux codes doivent donc être rattrapés et renvoyer la
        page d'accueil de l'application. Sans ce réglage, tout lien profond
        échoue au rafraîchissement — et les codes QR de réalité augmentée
        pointent précisément vers des liens profonds ;
  \item \textbf{les durées de cache différenciées} : les fichiers dont le nom
        contient une empreinte sont conservés un an, les modèles
        tridimensionnels une semaine, tandis que la page d'accueil et le
        travailleur de service ne sont jamais mis en cache — faute de quoi le
        site resterait figé sur une version périmée pendant plusieurs jours.
\end{itemize}

\begin{figure}[htbp]
\centering
\aplacer{9cm}{Représenter l'architecture d'hébergement en nuage effectivement
déployée. De gauche à droite : \textbf{Visiteur} (navigateur, téléphone)
$\rightarrow$ \textbf{Route 53} (zone hébergée, enregistrements A et AAAA pour
le domaine principal et enregistrements joker \texttt{*.} pour les
sous-domaines des organisations) $\rightarrow$ \textbf{CloudFront} (réseau de
diffusion, certificat \textbf{ACM}, politique d'en-têtes de sécurité, repli
applicatif sur les codes 403 et 404, durées de cache différenciées)
$\rightarrow$ \textbf{Amazon S3} (bucket privé, chiffré, versionné, accessible
uniquement par contrôle d'accès à l'origine). Ajouter en dérivation :
\textbf{Supabase} (base de données, authentification, fonctions serveur)
interrogé directement par le navigateur, et \textbf{EC2 à processeur graphique}
sollicité ponctuellement pour la photogrammétrie. Indiquer les 29 ressources
décrites par Terraform.}
\caption{Architecture de l'infrastructure d'hébergement en nuage}
\label{fig:cloud}
\source{Élaboré par nos soins d'après la description Terraform du projet.}
\end{figure}

\subsection{La chaîne d'intégration et de déploiement continus}

La chaîne mise en place comporte deux enchaînements.

Le premier, déclenché à chaque dépôt de modification, installe les dépendances,
compile l'application et exécute une série de contrôles portant sur le résultat
produit : présence des fichiers attendus, validité de l'en-tête des modèles
tridimensionnels, présence effective de l'adresse du service de données dans le
paquet compilé. Ce dernier contrôle mérite un mot : une compilation dépourvue de
cette adresse produit un site qui s'affiche parfaitement et ne charge aucune
donnée — la panne la plus déroutante qui soit. La description de l'infrastructure
est validée dans le même mouvement, sans aucun accès au fournisseur de nuage.

Le second, déclenché sur la branche principale, publie l'application en quatre
passes distinctes, chacune correspondant à une famille de fichiers ayant sa
durée de cache et son étiquette de type propres. Il purge ensuite le cache de
diffusion, puis — point essentiel — \textbf{vérifie le site en ligne} : code de
réponse de la page d'accueil, comportement du repli sur un lien profond, et
étiquette de type effectivement servie pour un modèle tridimensionnel. On ne
contrôle pas ce que l'on croit avoir déployé, on contrôle ce que le visiteur
reçoit.

L'authentification auprès du fournisseur de nuage ne repose sur aucune clé
permanente : un jeton de courte durée est présenté à chaque exécution, et le
rôle correspondant est restreint au dépôt de code concerné, à ce seul espace de
stockage et à cette seule distribution.

\begin{figure}[htbp]
\centering
\aplacer{8.5cm}{Représenter la chaîne DevOps actuellement en service. De gauche
à droite : \textbf{Poste de développement} (Git) $\rightarrow$ \textbf{Dépôt
GitHub} $\rightarrow$ \textbf{Intégration continue} (installation, compilation,
contrôles du paquet produit, \texttt{terraform fmt} et \texttt{validate})
$\rightarrow$ \textbf{Déploiement continu} sur la branche principale
(authentification par jeton OIDC de courte durée, publication en quatre passes
vers \textbf{S3}, purge du cache \textbf{CloudFront}, vérification a posteriori
du site en ligne). Faire figurer, en parallèle, la voie \textbf{Docker} : image
multi-étapes (compilation Node puis diffusion Nginx) utilisée pour reproduire en
local le comportement de production. Indiquer les points de contrôle et les
conditions d'arrêt de la chaîne en cas d'échec.}
\caption{Chaîne d'intégration et de déploiement continus en service}
\label{fig:devops-actuel}
\source{Élaboré par nos soins d'après les fichiers de description des chaînes du projet.}
\end{figure}

% =============================================================================
\section{Présentation des résultats : les fonctionnalités livrées}
\label{sec:resultats}

Cette section décrit chacune des fonctionnalités effectivement livrées et, pour
chacune, la finalité qu'elle poursuit. Le tableau~\ref{tab:fonctionnalites}
en donne d'abord une vue d'ensemble.

\begin{table}[htbp]
\centering
\caption{Vue d'ensemble des fonctionnalités livrées et de leur finalité}
\label{tab:fonctionnalites}
\footnotesize
\begin{tabularx}{\textwidth}{|L{1.0cm}|L{4.0cm}|X|C{1.6cm}|}
\hline
\rowcolor{keycetete}
\thead{N\textsuperscript{o}} & \thead{Fonctionnalité} & \thead{Finalité poursuivie} & \thead{Objectif visé} \\ \hline
F1 & Gestion des collections & Constituer un inventaire structuré et publiable des musées, salles et œuvres & OS1 \\ \hline
F2 & Site public de l'institution & Offrir une vitrine éditable par l'institution elle-même & OS1 \\ \hline
F3 & Visite immersive à 360 degrés & Restituer le déplacement dans les espaces & OS1 \\ \hline
F4 & Présentation en trois dimensions & Permettre l'examen d'une œuvre sous tous ses angles & OS1 \\ \hline
F5 & Réalité augmentée à taille réelle & Faire éprouver les dimensions véritables de l'objet & OS1 \\ \hline
F6 & Codes QR de passage à l'appareil mobile & Permettre le passage de l'ordinateur au téléphone sans installation & OS1 \\ \hline
F7 & Relief 2.5D & Révéler le volume d'une photographie plate lorsque le modèle 3D n'existe pas & OS1 \\ \hline
F8 & Assistant conversationnel ancré & Répondre aux questions sans inventer & OS3 \\ \hline
F9 & Guide vocal & Restituer le récit d'une œuvre à l'oral & OS3 \\ \hline
F10 & Journal des questions sans réponse & Signaler à l'institution les lacunes de sa documentation & OS3 \\ \hline
F11 & Mémoire réunifiée & Retrouver les œuvres apparentées dispersées dans le monde & OS2 \\ \hline
F12 & Cabinet de comparaison & Soumettre les rapprochements à la validation du conservateur & OS2 \\ \hline
F13 & Généalogie navigable & Relier une œuvre au chef qui l'a possédée et à sa lignée & OS2 \\ \hline
F14 & Billetterie à code QR & Vendre et contrôler les entrées & OS4 \\ \hline
F15 & Boutique en ligne & Ouvrir un canal de vente aux artisans & OS4 \\ \hline
F16 & Adhésions et campagnes & Transformer une visite en relation durable & OS4 \\ \hline
F17 & Cloisonnement multi-organisations & Accueillir plusieurs institutions sans mélange des données & OS4 \\ \hline
F18 & Assistant d'installation & Permettre à une institution de s'installer par conversation & OS4 \\ \hline
F19 & Statistiques de fréquentation & Mettre en avant automatiquement les œuvres regardées & OS4 \\ \hline
F20 & Carte de partage & Diffuser une œuvre sur les messageries & OS4 \\ \hline
\end{tabularx}
\source{Élaboré par nos soins.}
\end{table}

\subsection{F1 — La gestion des collections}

\textbf{Ce que fait la fonctionnalité.} Le back-office permet d'enregistrer les
musées, puis les salles qui les composent, puis les œuvres qu'abritent ces
salles. Chaque œuvre porte un nom, un nom commun, une description, une histoire,
des médias, une grille tarifaire, un indice de rareté et, le cas échéant, ses
modèles tridimensionnels.

\textbf{Pourquoi elle existe.} Elle constitue le socle de tout le reste : sans
inventaire structuré, ni la visite immersive, ni l'assistant, ni le
rapprochement documentaire ne disposent de matière. Elle répond directement à
l'absence de système de gestion relevée lors de l'étude de l'existant.

\textbf{Point de conception notable.} La publication obéit à une règle de
cascade : une œuvre n'apparaît au public que si son musée, sa salle \emph{et}
elle-même sont publiés. Cette règle évite qu'une pièce en cours de traitement
ne se retrouve exposée par inadvertance.

\begin{figure}[htbp]
\centering
\acapturer{6.5cm}{Insérer la capture d'écran de la liste des œuvres dans le
back-office, montrant la barre d'outils, la recherche, les vignettes, les
indicateurs de publication et les boutons d'action de ligne.}
\caption{Back-office : écran de gestion des œuvres}
\label{fig:cap-objets}
\source{Capture d'écran de la plateforme MUSÉA, août 2026.}
\end{figure}

\subsection{F2 — Le site public de l'institution}

\textbf{Ce que fait la fonctionnalité.} Chaque institution dispose d'un site
public présentant son identité, ses musées, son catalogue, ses événements et sa
boutique.

\textbf{Pourquoi elle existe.} Pour que l'institution reste maîtresse de son
image. L'ensemble des éléments visibles — marque, couleurs, images d'accueil,
textes des pages de connexion, blocs de la page d'accueil — est modifiable
depuis le back-office. Aucun texte n'est figé dans le code : une institution qui
change de charte graphique n'a pas besoin d'un développeur.

\subsection{F3 — La visite immersive à 360 degrés}

\textbf{Ce que fait la fonctionnalité.} Le visiteur entre dans une salle,
pivote librement, zoome, passe à la salle suivante par un point de passage, et
touche un objet pour en ouvrir la fiche. Une mini-carte et une barre de
progression lui indiquent où il se trouve dans le parcours.

\textbf{Pourquoi elle existe.} Elle répond à l'attente la plus fréquemment
exprimée par les visiteurs interrogés — se déplacer dans les salles comme si
l'on y était — et à la première limite relevée dans la problématique : la visite
passive sans sensation d'espace.

\textbf{Point de conception notable.} Un mode de démonstration dessine une salle
stylisée au vol, avec plafond, cimaise, alcôves éclairées et sol, dans une
teinte différente par salle. Le parcours est donc démontrable \emph{avant} les
prises de vue réelles — ce qui, sur un terrain dépourvu de matière numérique
préexistante, a conditionné la possibilité même de présenter le dispositif.

\begin{figure}[htbp]
\centering
\acapturer{6.5cm}{Insérer la capture d'écran du parcours immersif côté visiteur :
panorama à 360 degrés, points d'intérêt cliquables, mini-carte, barre de
progression et fiche flottante d'une œuvre.}
\caption{Site public : le parcours de visite immersive}
\label{fig:cap-tour}
\source{Capture d'écran de la plateforme MUSÉA, août 2026.}
\end{figure}

\subsection{F4 et F5 — La présentation en trois dimensions et la réalité augmentée}

\textbf{Ce que font les fonctionnalités.} Depuis la fiche d'une œuvre, le
visiteur fait pivoter le modèle tridimensionnel, l'examine sous tous ses angles,
puis — s'il utilise un téléphone compatible — le pose dans son propre
environnement, à ses dimensions véritables.

\textbf{Pourquoi elles existent.} Pour répondre à la deuxième attente des
visiteurs et pour donner corps à ce que nous avons appelé le « retour au pays » :
faire apparaître, dans le salon d'un membre de la diaspora, un objet conservé à
Bangoulap, à sa taille réelle. Une photographie ne produit jamais cet effet.

\textbf{Couverture des appareils.} Le tableau~\ref{tab:appareils} présente
l'état de la couverture constatée.

\begin{table}[htbp]
\centering
\caption{Couverture des appareils pour la réalité augmentée}
\label{tab:appareils}
\small
\begin{tabularx}{\textwidth}{|L{5.4cm}|C{2.6cm}|X|}
\hline
\rowcolor{keycetete}
\thead{Plateforme} & \thead{État} & \thead{Condition} \\ \hline
Android avec ARCore & Fonctionnel & Aucune ; y compris sur la pièce de démonstration \\ \hline
iPhone et iPad & Fonctionnel & Un modèle au format USDZ doit être chargé sur l'œuvre \\ \hline
Android sans interface WebXR & Fonctionnel & Les modèles doivent être des fichiers statiques et non des adresses temporaires \\ \hline
Ordinateur & Non applicable & Un code QR renvoie vers le téléphone \\ \hline
\end{tabularx}
\source{Essais conduits en août 2026.}
\end{table}

\subsection{F6 — Les codes QR de passage à l'appareil mobile}

\textbf{Ce que fait la fonctionnalité.} Un code QR affiché à l'écran renvoie
vers une page autonome dédiée à la réalité augmentée pour l'œuvre concernée.

\textbf{Pourquoi elle existe.} Parce que la réalité augmentée exige un téléphone
tandis qu'une présentation se tient devant un ordinateur. Sans ce pont, la
fonctionnalité la plus démonstrative du projet resterait invisible au public
auquel on la présente.

\subsection{F7 — Le relief 2.5D}

\textbf{Ce que fait la fonctionnalité.} Une photographie plate révèle du volume
lorsque l'utilisateur bouge : les sommets d'un maillage subdivisé sont déplacés
selon une carte de profondeur estimée.

\textbf{Pourquoi elle existe.} Toutes les œuvres ne seront jamais numérisées en
trois dimensions : l'opération est longue et suppose de manipuler la pièce. Le
relief offre un intermédiaire entre la photographie plate et le modèle complet.

\textbf{Limite assumée.} L'illusion tient de face et s'effondre de profil.
L'angle de parallaxe est donc borné à vingt degrés, au-delà desquels les zones
absentes de la photographie apparaissent et s'étirent. Ce bridage est la
condition du procédé, non un réglage de confort.

\subsection{F8 — L'assistant conversationnel à ancrage documentaire}

\textbf{Ce que fait la fonctionnalité.} Le visiteur pose une question en
français ou en anglais ; l'assistant répond à partir des documents de
l'institution, cite ses sources et propose des liens vers les fiches concernées.

\textbf{Pourquoi elle existe.} Pour répondre à la troisième limite de la
problématique — l'absence d'aide face aux questions précises — sans encourir le
risque que redoutait le personnel : qu'un dispositif automatique énonce des
contrevérités sur le patrimoine.

\textbf{Garde-fous mis en place.} Quatre mécanismes concourent à la fiabilité :
un refus poli des sujets hors périmètre, prononcé sans même solliciter le modèle
de langue ; une réponse chaleureuse aux salutations, plutôt qu'une demande de
précision déplacée ; une restriction de la recherche au musée ou à la salle où
se trouve le visiteur ; et une consigne interdisant d'énoncer un fait absent des
documents fournis, assortie de l'autorisation explicite de déclarer son
ignorance.

\begin{figure}[htbp]
\centering
\acapturer{6.5cm}{Insérer la capture d'écran de l'assistant conversationnel :
question du visiteur, réponse ancrée avec mention des sources, et cartes de
liens vers les fiches d'œuvres concernées.}
\caption{Site public : l'assistant conversationnel en fonctionnement}
\label{fig:cap-guide}
\source{Capture d'écran de la plateforme MUSÉA, août 2026.}
\end{figure}

\subsection{F9 — Le guide vocal}

\textbf{Ce que fait la fonctionnalité.} Le récit d'une œuvre ou d'une salle peut
être écouté, avec une voix choisie par l'institution et pouvant varier d'un
musée à l'autre, voire d'une salle à l'autre.

\textbf{Pourquoi elle existe.} Parce que la médiation muséale est d'abord orale
et parce que l'écoute libère le regard : on ne peut pas lire une notice et
observer une œuvre en même temps.

\subsection{F10 — Le journal des questions restées sans réponse}

\textbf{Ce que fait la fonctionnalité.} Les questions auxquelles l'assistant n'a
pu répondre sont consignées de manière anonyme, classées en huit thèmes
— matière, datation, origine, usage, auteur, dimensions, valeur, conservation —
et présentées au personnel sous forme de conseils.

\textbf{Pourquoi elle existe.} Elle retourne le dispositif au profit de
l'institution : au lieu de constater un échec, elle indique quoi documenter.
Un conseil tel que « précisez la matière et la technique de fabrication » est
directement actionnable, là où un simple décompte ne le serait pas.

\textbf{Point de conception notable.} La classification n'emploie aucun modèle
de langue. Ce choix est délibéré : huit catégories fermées produisent un
affichage instantané et, surtout, continuent de fonctionner lorsque aucune clé
d'accès n'est disponible — situation courante dans le contexte du projet.

\subsection{F11 et F12 — La mémoire réunifiée et le cabinet de comparaison}

\textbf{Ce que font les fonctionnalités.} À partir de la notice d'une œuvre, le
système interroge en parallèle cinq collections ouvertes, ordonne les résultats
par proximité de sens, fait justifier les rapprochements les plus probables,
puis les soumet au conservateur qui valide, écarte ou corrige la justification.
Les correspondances validées apparaissent alors sur la fiche publique de
l'œuvre, avec le nom du musée détenteur et le numéro d'inventaire.

\textbf{Pourquoi elles existent.} Pour répondre à la deuxième limite de la
problématique — des objets isolés et une histoire incomplète — et au thème
inattendu apparu lors des entretiens.

\textbf{Résultat marquant.} Une mesure comparative conduite sur les mêmes
collections et au même seuil de notation établit ce qui suit : une recherche
formulée directement en français ramène 16 résultats dont le meilleur atteint un
score de 57 sur 100, tandis que la même recherche reformulée dans le vocabulaire
de catalogage anglais n'en ramène que 5, mais dont le meilleur atteint 85. Ce
résultat est développé et interprété au chapitre~4.

\textbf{Point de conception notable.} Le catalogue des objets externes est
mondial et non rattaché à une organisation : un masque conservé à New York est
le même objet pour toutes les institutions de la plateforme. Il n'est donc
interrogé et vectorisé qu'une seule fois, ce qui fait \emph{décroître} le coût
par institution à mesure que la plateforme s'étoffe.

\begin{figure}[htbp]
\centering
\acapturer{6.5cm}{Insérer la capture d'écran du cabinet de comparaison côté
public : l'œuvre de la Fondation au centre, les œuvres apparentées disposées en
arc, chacune accompagnée de son cartel — musée détenteur, numéro d'inventaire et
justification du rapprochement.}
\caption{Site public : le cabinet de comparaison des œuvres apparentées}
\label{fig:cap-freres}
\source{Capture d'écran de la plateforme MUSÉA, août 2026.}
\end{figure}

\subsection{F13 — La généalogie navigable}

\textbf{Ce que fait la fonctionnalité.} Un arbre bidirectionnel présente les
ascendants à gauche, les descendants à droite et la personne consultée au
centre. Un clic recentre l'arbre sur une autre personne. Une recherche
instantanée met en évidence les cartes correspondantes. Le chemin de filiation
entre deux personnes — par l'aïeul commun le plus proche — est calculé et mis en
évidence. Une frise chronologique présente enfin la succession des règnes.

\textbf{Pourquoi elle existe.} Pour donner corps au fil narratif qui structure
tout le projet : \emph{objet $\rightarrow$ chef qui l'a possédé $\rightarrow$
chefferie $\rightarrow$ lignée $\rightarrow$ autres objets de la lignée}. Un
objet qui ne mène nulle part est un échec de conception.

\textbf{Point de conception notable.} Les fins de règne non renseignées sont
estimées à l'avènement du successeur, mais signalées par des hachures plutôt que
présentées comme des faits établis. De même, répondre « aucun lien connu » est
traité comme une réponse assumée et non comme un échec.

\subsection{F14 et F15 — La billetterie et la boutique}

\textbf{Ce que font les fonctionnalités.} Le visiteur achète un billet ou un
produit d'artisanat ; le billet est délivré sous forme de code QR ; le personnel
le contrôle à l'entrée au moyen de la caméra d'un téléphone.

\textbf{Pourquoi elles existent.} Pour ouvrir un canal de recette qui n'existait
pas et pour remplacer le registre papier par un contrôle vérifiable. La boutique
répond en outre à une préoccupation exprimée par la Fondation : donner un
débouché à la production des ateliers indépendamment des visites physiques.

\subsection{F16 — Les adhésions et les campagnes d'information}

\textbf{Ce que fait la fonctionnalité.} Un visiteur qui crée un compte sur le
site d'une institution en devient l'adhérent. L'institution dispose alors d'un
fichier exportable et peut lui adresser des campagnes d'information, dont la
rédaction peut être assistée.

\textbf{Pourquoi elle existe.} Pour transformer une visite ponctuelle en
relation durable, ce que l'absence de fichier de visiteurs interdisait
totalement.

\textbf{Garanties apportées.} Les destinataires sont calculés par la base de
données et restreints aux adhérents ayant consenti à être contactés ; chaque
message comporte un lien de désinscription ; le contenu saisi par
l'administrateur est intégralement échappé, de sorte qu'aucun code ne puisse y
être injecté ; et la rédaction assistée se voit interdire d'inventer une date,
un horaire, un tarif ou un nom d'œuvre qui ne lui aurait pas été fourni.

\subsection{F17 — Le cloisonnement multi-organisations}

\textbf{Ce que fait la fonctionnalité.} Chaque institution vit sur son propre
sous-domaine et ne voit que ses propres données. Trois cercles d'identité
coexistent sans jamais se confondre : le super-administrateur de la plateforme,
le personnel d'une institution, et le visiteur — lequel possède un compte global
et adhère à une ou plusieurs institutions.

\textbf{Pourquoi elle existe.} Pour répondre au quatrième objectif : permettre
l'accueil de nouvelles institutions sans reconstruction. C'est la condition de
la diffusion du dispositif au-delà de la Fondation.

\textbf{Point de conception notable.} Le cloisonnement est appliqué par la base
de données au moyen de règles de sécurité au niveau de la ligne, et non par le
code de l'application. La différence est décisive : même un composant détourné
ou défectueux ne peut pas lire les données d'une autre institution, parce que
c'est la base qui refuse, et non le programme qui s'abstient de demander.

\subsection{F18 — L'assistant d'installation}

\textbf{Ce que fait la fonctionnalité.} Une institution qui vient de s'inscrire
décrit ses collections en langage courant ; un agent crée pour elle les musées,
les salles et les premières œuvres.

\textbf{Pourquoi elle existe.} Parce que la page blanche est le premier obstacle
à l'adoption. Une institution qui doit saisir cent notices avant de voir quoi
que ce soit abandonne.

\textbf{Garantie de sécurité.} L'agent écrit avec le jeton de la personne qui
l'utilise, jamais avec une clé privilégiée. Les règles de cloisonnement
s'appliquent donc exactement comme si la personne cliquait elle-même : un agent
détourné par une consigne malveillante ne peut rien écrire hors de son
institution. Les outils dont il dispose sont en outre limités à une liste
blanche — création de musée, de salle, d'œuvre et d'identité — sans aucune
suppression, et le nombre de créations est plafonné.

\subsection{F19 et F20 — Les statistiques de fréquentation et la carte de partage}

\textbf{Ce que font les fonctionnalités.} La première comptabilise les
consultations des fiches d'œuvres et met automatiquement en avant, sur la page
d'accueil, les pièces les plus regardées. La seconde produit une image carrée
prête à être partagée sur les messageries, comportant la photographie de
l'œuvre, son nom, sa salle, la marque de l'institution et un code QR renvoyant
vers sa fiche.

\textbf{Pourquoi elles existent.} L'une et l'autre travaillent sans que le
personnel ait à saisir, à consulter ou à décider quoi que ce soit — critère
retenu pour juger de leur utilité réelle dans une institution aux effectifs
restreints.

\textbf{Point de conception notable.} La mise en avant automatique n'intervient
qu'au-delà d'un seuil de consultations ; en deçà, la sélection humaine est
conservée, car trois clics ne font pas une tendance.

% =============================================================================
\section{Synthèse des mesures obtenues}

Le tableau~\ref{tab:mesures} rassemble les mesures relevées sur le système, qui
serviront de base à la vérification des hypothèses au chapitre suivant.

\begin{table}[htbp]
\centering
\caption{Synthèse des mesures relevées sur le système}
\label{tab:mesures}
\small
\begin{tabularx}{\textwidth}{|C{1.7cm}|X|C{3.0cm}|}
\hline
\rowcolor{keycetete}
\thead{Variable} & \thead{Indicateur} & \thead{Valeur relevée} \\ \hline
V1 & Dimensions mesurées sur le modèle de démonstration & 45 $\times$ 52 $\times$ 45 cm \\ \hline
V1 & Sommets du maillage de démonstration & 835 \\ \hline
V1 & Plateformes couvertes pour la réalité augmentée & 3 sur 3 \\ \hline
V1 & Angle de parallaxe admissible pour le relief & $\pm$ 20\textdegree \\ \hline
V2 & Collections ouvertes effectivement interrogées & 5 \\ \hline
V2 & Œuvres apparentées identifiées & 41 \\ \hline
V2 & Meilleur score de correspondance obtenu & 90 / 100 \\ \hline
V2 & Pays représentés sur le globe de dispersion & 137 \\ \hline
V2 & Score en recherche française brute (meilleur) & 57 / 100 \\ \hline
V2 & Score après reformulation en vocabulaire de catalogue & 85 / 100 \\ \hline
V2 & Séparation des scores en anglais (frères / bruit) & 0,898 / 0,779 \\ \hline
V2 & Séparation des scores en français (étendue) & 0,13 \\ \hline
V3 & Fonctions serveur déployées & 12 \\ \hline
V3 & Fragments vectorisables par appel & 18 (lots de 12 retenus) \\ \hline
V3 & Durée de vectorisation de 18 fragments & 1,7 s \\ \hline
V3 & Thèmes de classement des questions sans réponse & 8 \\ \hline
V4 & Tables de la base de données & 38 \\ \hline
V4 & Ressources d'infrastructure décrites par le code & 29 \\ \hline
V4 & Ressources effectivement créées à la mise en ligne & 19 \\ \hline
V4 & Coût récurrent mensuel de l'hébergement & 730 FCFA \\ \hline
\end{tabularx}
\source{Mesures effectuées sur la plateforme, juillet--août 2026.}
\end{table}

\conclchap

Ce chapitre a présenté le terrain — la Fondation Jean Félicien Gacha et son
\emph{Think Tank} Foudjem —, le déroulement du stage, les données recueillies
auprès du personnel et des visiteurs, puis le cahier des charges, la
modélisation, l'implémentation et les vingt fonctionnalités livrées. Les
mesures rassemblées en fin de chapitre constituent la matière factuelle du
diagnostic. Le chapitre suivant les interprète, vérifie les hypothèses formulées
en introduction et présente l'intervention proposée pour la suite.
